\begin{lemma}\label{lem:decic48cusp13}
Let $X=X_1(48)_{\mathbf F_5}$, and put
$A_{13}=T_{13}-1-13\langle13\rangle$.
For every $\mathbf F_5$-rational cuspidal divisor $C$ on $X$, the divisor
$A_{13}C$ is principal.
\end{lemma}
\begin{proof}
We use primitive cusp labels $(a,b)\in(\mathbf Z/48\mathbf Z)^2$, modulo
$(a,b)\sim\pm(a+nb,b)$. The arithmetic and diamond actions are
\[
 \sigma_u(a,b)=(ua,b),\qquad
 \langle u\rangle(a,b)=(ua,u^{-1}b).
\]
There are $56$ geometric cusps. Their $\sigma_5$-orbits consist of
$16$ rational cusps, $12$ quadratic cusps, and $4$ quartic cusps.
On cusps one has
\[
 T_{13}=\sigma_{13}+13\langle13\rangle\sigma_{13}^{-1}.
\]
Applying this identity to the $32$ closed cusps shows that $A_{13}$ vanishes
on all but four of them. These four form two pairs: the two quadratic
cusps of width $2$, and the two quartic cusps of width $1$.
If their differences are denoted by $C_2-C_2'$ and $C_4-C_4'$, respectively,
the nonzero columns of $A_{13}$ are
\[
 \pm12(C_2-C_2'),\qquad \pm12(C_4-C_4').
\]
This calculation is a permutation calculation on primitive labels;
\path{independent_cusp_A13_labels.json} records every column.

Both displayed differences, multiplied by $12$, are principal. More
precisely, the certificate \path{cusp_corrections_v2.m} gives the two
prime pairs and rational functions $f_2,f_4$ in the degree-$19$
function-field presentation of $X$, with
\[
 \operatorname{div}(f_i)=-12(C_i-C_i')\qquad(i=2,4).
\]
The primes are characterized by their degrees and their pole orders
$2$ and $1$ for $j$, respectively. The identities were checked in both
the finite and infinite maximal orders of a fresh copy of the function
field; see \path{independent_cusp_principal.m}. This proves the claim.
\end{proof}

\begin{lemma}\label{lem:decic48NC13}
Let $D$ be an effective rational divisor of degree $10$ on $X_1(48)$,
and write its reduction at $5$ as $\overline D=N+C$, with $C$ cuspidal
and $N$ supported away from the cusps. Then $A_{13}N$ is principal.
If $C=0$, the same conclusion holds with $13$ replaced by $7$.
\end{lemma}
\begin{proof}
The group $J_1(48)(\mathbf Q)$ is finite by DEHMZ, Theorem~3.1(2).
At either good odd prime $\ell=7,13$, the Eichler--Shimura relation and
injectivity of reduction on rational torsion give
\[
 (T_\ell-1-\ell\langle\ell\rangle)J_1(48)(\mathbf Q)=0.
\]
The same operator annihilates the degree part after subtracting a rational
cusp. Such a cusp is annihilated as a divisor. Reducing at $5$ therefore
gives $A_{13}(N+C)\sim0$. Lemma~\ref{lem:decic48cusp13} gives
$A_{13}C\sim0$, so $A_{13}N\sim0$. If $C=0$, the argument applies
directly to $A_7$.
\end{proof}
